\documentclass{situationist-preprint}

\title{The Bald Man Paradox}
\msid{LNG-12-c}
\msfamily{LNG-12}
\msdiscipline{Philosophy of Language}
\mstype{Consequence}
\mskeywords{vagueness; sorites paradox; baldness; gradable adjectives; epistemicism; multidimensional vagueness}

\begin{document}
\maketitle

\begin{abstract}
The bald man paradox, the twin of the heap in the ancient sources, runs in the opposite direction and concerns a predicate that differs from \enquote{heap} in three ways: whether a man is bald depends on how his hair is distributed and not only on how many hairs he has; \enquote{bald} is a gradable adjective whose standard varies with context; and a man becomes bald over time. We apply the account of LNG-12-a to each feature. The first yields a sharp boundary along every series of cases from a clear non-bald to a clear bald head, though not a single number of hairs. The second relativizes the boundary to a standard without making it vague. The third yields a last hair whose loss made a man bald, which no one can identify.
\end{abstract}

\section{The paradox and the account}

Diogenes Laertius reports that Eubulides of Miletus propounded, among other puzzles, the heap and the bald man \citep[II.108]{DiogenesLaertius}. In its standard form: a man with a full head of hair is not bald; removing a single hair never makes a man who is not bald bald; so a man with no hairs is not bald. LNG-12-a argues, for \enquote{heap}, that such an argument is valid, that its tolerance premise is false, and that the boundary its falsity implies cannot be known because the meaning of the word is plastic. Its argument uses three premises about the predicate: bivalence under every pattern of use (Assumption 2.2), clear cases at the ends of the series (Assumption 2.3), and monotonicity along it (Assumption 2.4). The bald man differs from the heap in ways that bear on the third of these and on what the boundary is a boundary of. We take the differences in turn.

\section{Distribution: boundaries along series}

Whether a man is bald depends on where his hairs are. A man with many hairs at the back and sides and none on the crown is bald; a man with fewer hairs evenly spread may not be. So there is no number $n$ such that every man with $n$ hairs is bald and every man with $n+1$ is not, and monotonicity in the number of hairs fails. The cases of baldness are ordered not linearly but partially: one head is at least as bald as another if it can be obtained from it by removing hairs. Most vague predicates are multidimensional in this way \citep{Keefe2000}, and the objection to sharp boundaries is that a partial order leaves no single place for a boundary to be.

The account needs no single place. Call a \emph{series} a finite sequence of heads $h_0, \dots, h_m$ in which each is obtained from its predecessor by removing one hair, $h_0$ is clearly not bald and $h_m$ is clearly bald. Monotonicity holds along every series, since removing hairs never makes a man less bald.

\begin{proposition}\label{p:series}
Under every pattern of use that could be ours, every series contains exactly one $i$ such that $h_i$ is not bald and $h_{i+1}$ is bald.
\end{proposition}

\begin{proof}
Along a series, \enquote{bald} satisfies the three premises that Lemma 3.1 and Theorem 3.2 of LNG-12-a use for \enquote{heap}, with heads indexed by position in place of numbers of grains. The proof of those results uses nothing about grains beyond their linear order, so it applies.
\end{proof}

The boundary of \enquote{bald} under a pattern of use is then not a number but a set of heads: the maximal non-bald heads, those that are not bald while every head obtained from them by removing a hair is bald. Every series crosses it exactly once. This is what the sharpness of a multidimensional predicate consists in, and it is no less sharp for being a frontier rather than a point. Its unknowability follows as in Theorem 3.5 of LNG-12-a: for each head on the frontier, a pattern of use too close to ours to be discriminated from it places that head just inside or just outside.

\section{Context: boundaries relative to standards}

\enquote{Bald} is gradable: one man is balder than another, and \enquote{completely bald} and \enquote{rather bald} modify it as they modify other gradable adjectives. LNG-19 holds that gradable adjectives are context-sensitive, in that context supplies a standard on the scale they introduce \citep{Kennedy2007}. A barber and a casting director may apply \enquote{bald} differently without either misusing it.

Context-sensitivity is not vagueness, and it does not threaten the result. Fix a context, and with it a standard. Relative to that standard, the premises of Proposition~\ref{p:series} hold, and so there is a sharp frontier, unknowable by those whose use fixes it. Varying the context changes which frontier is in play; it does not make any frontier blurred. The sorites does not exploit context: it is compelling within a fixed context, as when a single barber is asked, hair by hair, whether his client is now bald. Accounts that locate the source of tolerance in shifts of context during a forced march, such as that of \citet{Raffman1994}, must explain why the premise seems true in a context held fixed; the epistemic account explains it by the margin for error discussed in LNG-12-b, which holds within every context.

\section{Time: the hair that made him bald}

A man who loses his hair one by one traces a series through time. By Proposition~\ref{p:series}, under our pattern of use there was a last hair whose loss made him bald: before it fell he was not bald, after it fell he was. It sounds absurd to say that a particular hair, falling on a particular morning, made the difference.

\begin{objection}
If a single hair made the difference, its loss was a significant event in the man's life. But nothing about that morning was significant. So no hair made the difference.
\end{objection}

\begin{reply}
The loss changed which of our words apply to him, not anything that matters to him or to anyone. The event was significant for the extension of \enquote{bald} and for nothing else, and since no one can know when it occurred, no one could treat it as significant. The appearance of absurdity comes from conflating the sharpness of a boundary with the importance of crossing it. That a sharp boundary in a classification marks no important change in the classified is familiar from legal thresholds, and the epistemic account adds only that for vague terms the threshold is drawn by use rather than by decree.
\end{reply}

\begin{objection}
The epistemic account must say that the man's barber, who saw him every week, did not know whether he was bald for many weeks. But some of those weeks the barber would have answered with confidence.
\end{objection}

\begin{reply}
Confidence in a borderline case is consistent with ignorance. A barber who confidently says \enquote{bald} of a head just over the frontier is right, but by luck: a pattern of use he could not have told apart from ours would have made him wrong. On the account, that is why we do not credit borderline verdicts as knowledge, however firmly delivered \citep[ch.~8]{Williamson1994}.
\end{reply}

\section{Conclusion}

The bald man adds nothing to the heap that the account of LNG-12-a cannot absorb. Multidimensionality turns the boundary into a frontier crossed once by every series; gradability makes the frontier relative to a standard; and the passage of time makes the crossing an event. In each case the frontier is sharp, fixed by use, and unknowable, as MET-18 holds of vague predicates in general.

\begin{thebibliography}{9}\small
\bibitem[Diogenes Laertius(1925)]{DiogenesLaertius} Diogenes Laertius (1925). \emph{Lives of Eminent Philosophers}. Trans. R.~D. Hicks. Loeb Classical Library, Harvard University Press. Cited by book and section.
\bibitem[Keefe(2000)]{Keefe2000} Keefe, R. (2000). \emph{Theories of Vagueness}. Cambridge University Press.
\bibitem[Kennedy(2007)]{Kennedy2007} Kennedy, C. (2007). Vagueness and grammar: The semantics of relative and absolute gradable adjectives. \emph{Linguistics and Philosophy}, 30(1), 1--45.
\bibitem[Raffman(1994)]{Raffman1994} Raffman, D. (1994). Vagueness without paradox. \emph{The Philosophical Review}, 103(1), 41--74.
\bibitem[Williamson(1994)]{Williamson1994} Williamson, T. (1994). \emph{Vagueness}. Routledge.
\end{thebibliography}

\end{document}
